\section{Additional experimental details}

In this appendix, we provide additional experimental details.
Section~\ref{sec:regulariser_details} provides additional details about the regularisers we used and Sec.~\ref{sec:hyperparameters} details the training hyperparamters used for our experiments.

\begin{table*}[tb]
\caption{Training hyperparamters for experiments in the main paper.  ``--'' indicates that the regularisation method was not used at all. 
Values which are constant across all columns are listed once.
Datasets are denoted as follows: K400: Kinetics 400. K600: Kinetics 600. MiT: Moments in Time. EK: Epic Kitchens. SSv2: Something-Something v2.}
\begin{tabularx}{\linewidth}{lYYYYY}
\toprule
                              & K400 & K600 & MiT & EK & SSv2 \\ \midrule
\multicolumn{6}{l}{\textit{Optimisation}}                                                                                 \\
Optimiser               & \multicolumn{5}{c}{Synchronous SGD}                                                               \\
Momentum					& \multicolumn{5}{c}{0.9} \\
Batch size 						& \multicolumn{5}{c}{64} \\
Learning rate schedule	& \multicolumn{5}{c}{cosine with linear warmup} \\
Linear warmup epochs	& \multicolumn{5}{c}{2.5} \\
Base learning rate			& 0.1 & 0.1 & 0.25 & 0.5 & 0.5 \\
Epochs	& 30 & 30 & 10 & 50 & 35 \\
\midrule
\multicolumn{6}{l}{\textit{Data augmentation}} \\
Random crop probability & \multicolumn{5}{c}{1.0} \\
Random flip probability & \multicolumn{5}{c}{0.5} \\
Scale jitter probability	& \multicolumn{5}{c}{1.0} \\
Maximum scale			  & \multicolumn{5}{c}{1.33} \\
Minimum scale			  & \multicolumn{5}{c}{0.9} \\
Colour jitter probability  & 0.8 & 0.8 & 0.8 & -- & -- \\
Rand augment number of layers~\cite{cubuk_arxiv_2019}		 & -- & -- & -- & 2 & 2 \\
Rand augment magnitude~\cite{cubuk_arxiv_2019} & -- & -- & -- & 15 & 20 \\
\midrule
\multicolumn{6}{l}{\textit{Other regularisation}} \\
Stochastic droplayer rate, $p_{\text{drop}}$~\cite{huang_stochasticdepth_eccv_2016} & -- & -- & -- & 0.2 &  0.3 \\
Label smoothing $\lambda$~\cite{szegedy_cvpr_2016}				& -- & -- & -- & 0.2 & 0.3 \\
Mixup $\alpha$~\cite{zhang_mixup_iclr_2018}				 & -- & -- & -- & 0.1 & 0.3 \\
\bottomrule
\end{tabularx}
\label{tab:training_hyperparameters}
\end{table*}

\subsection{Further details about regularisers}
\label{sec:regulariser_details} 

In this section, we provide additional details and list the hyperparameters of the additional regularisers that we employed in Tab.~\ref{tab:regularisation_epic_kitchens}.
Hyperparameter values for all our experiments are listed in Tab.~\ref{tab:training_hyperparameters}.

\paragraph{Stochastic depth}

Stochastic depth regularisation was originally proposed for training very deep residual networks~\cite{huang_stochasticdepth_eccv_2016}.
Intuitively, the outputs of a layer, $\ell$, are ``dropped out'' with probability, $p_{\text{drop}}(\ell)$ during training, by setting the output of the layer to be equal to its input.

Following~\cite{huang_stochasticdepth_eccv_2016}, we linearly increase the probability of dropping a layer according to its depth within the network, 
\begin{equation}
p_{\text{drop}}(\ell) = \frac{\ell}{L} p_{\text{drop}},
\end{equation}
where $\ell$ is the index of the layer in the network, and $L$ is the total number of layers.

\paragraph{Random augment}
Random augment~\cite{cubuk_arxiv_2019} randomly applies data augmentation transformations sequentially to an input example.
We follow the public implementation\footnote{\url{https://github.com/tensorflow/models/blob/master/official/vision/beta/ops/augment.py}}, but modify the data augmentation operations to be temporally consistent throughout the video (in other words, the same transformation is applied on each frame of the video).

The authors define two hyperparameters for Random augment, ``number of layers'' , the number of augmentation transformations to apply sequentially to a video and ``magnitude'',  the strength of the transformation that is shared across all augmentation operations.
Our values for these parameters are shown in Tab.~\ref{tab:training_hyperparameters}.

\paragraph{Label smoothing}
Label smoothing was proposed by~\cite{szegedy_cvpr_2016} originally to regularise training Inception-v3.
Concretely, the label distribution used during training, $\tilde{y}$, is a mixture of the one-hot ground-truth label, $y$,  and a uniform distribution, $u$, to encourage the network to produce less confident predictions during training:
\begin{equation}
\tilde{y} = (1 - \lambda) y + \lambda u.
\end{equation}
There is therefore one scalar hyperparamter, $\lambda \in [0, 1]$.

\paragraph{Mixup}
Mixup~\cite{zhang_mixup_iclr_2018} constructs virtual training examples which are a convex combination of pairs of training examples and their labels.
Concretely, given $(x_i, y_i)$ and $(x_j, y_j)$ where $x_i$ denotes an input vector and $y_i$ a one-hot input label, mixup constructs the virtual training example,
\begin{align}
\tilde{x} &= \lambda x_i + (1 - \lambda) x_j \nonumber \\
\tilde{y} &= \lambda y_i + (1 - \lambda) y_j.
\end{align}
$\lambda \in [0, 1]$, and is sampled from a Beta distribution, $\text{Beta}(\alpha, \alpha)$.
Our choice of the hyperparameter $\alpha$ is detailed in Tab.~\ref{tab:training_hyperparameters}.


\subsection{Training hyperparameters}
\label{sec:hyperparameters}

Table~\ref{tab:training_hyperparameters} details the hyperparamters for all of our experiments.
We use synchronous SGD with momentum, a cosine learning rate schedule with linear warmup, and a batch size of 64 for all experiments.
As aforementioned, we only employed additional regularisation when training on the smaller Epic Kitchens and Something-Something v2 datasets.